% ID: SEL-D0063-OA-E01
% ESTADO: validada. S063.
Sean $x,y$ los números de contenedores de A y B. Hay que minimizar $C=350x+550y$ euros bajo las restricciones
\[
2x+y\ge50,\qquad 3x+5y\ge180,\qquad x+y\le50,\qquad x,y\ge0.
\]
Las dos cotas de demanda se cruzan en $(10,30)$. Al cruzar la cota de almacenamiento con cada una, se obtienen $(0,50)$ y $(35,15)$, respectivamente. Los tres puntos son factibles y delimitan el triángulo.
\begin{center}
\setlength{\unitlength}{0.85mm}
\begin{picture}(92,65)(-12,-7)
\qbezier(0,0)(40,0)(80,0)
\qbezier(0,0)(0,27)(0,54)
\qbezier(0,50)(10,40)(20,30)
\qbezier(20,30)(45,22.5)(70,15)
\qbezier(70,15)(35,32.5)(0,50)
\put(81,0){$x$}\put(2,54){$y$}
\put(-12,48){\scriptsize $(0,50)$}
\put(15,23){\scriptsize $(10,30)$}
\put(70,10){\scriptsize $(35,15)$}
\end{picture}
\end{center}
Los costes son $20\,000$, $27\,500$ y $20\,500$ euros en $(10,30)$, $(0,50)$ y $(35,15)$, respectivamente. Una función lineal alcanza el mínimo en algún vértice. Como el mejor vértice ya es entero, también resuelve el problema de contenedores completos. El pedido óptimo es de $10$ a A y $30$ a B, por $\boxed{20\,000\text{ euros}}$.
